\chapter{CONCLUSÃO}\label{sec:conclusao}

Este trabalho comparou duas formas de obter o ângulo de deriva com o sensor virtual MGGP da IC: como terceira saída do modelo (28 modelos 5×3 aprovados) ou calculado a partir das velocidades estimadas (30 modelos 5×2), em simulação livre em três pistas de teste.

Quanto à primeira pergunta, a resposta depende da pista. Em duas das três pistas de teste não houve diferença significativa, e os intervalos de confiança não permitem excluir diferenças de até cerca de 0,2°. Na Wang 3.2, os modelos 5×3 reduziram a mediana do RMSE de $\beta$ de 1,74° para 1,10°, com efeito grande. Esse resultado vem de uma única pista e, nessa pista, como nas demais de teste, a velocidade lateral de referência está deslocada em cerca de 0,3~m/s em relação ao valor implícito no ângulo de deriva bruto, o que pode favorecer ou prejudicar um braço conforme o viés de cada modelo. De fato, a reavaliação com o deslocamento removido inverte a comparação: com a referência corrigida, os modelos 5×2 erram menos do que os 5×3 nas três pistas de teste (diferenças de medianas de $+0{,}16^\circ$ a $+0{,}45^\circ$). Como os dados não permitem dizer qual das referências está correta, o ganho do 5×3 na Wang 3.2 não pode ser tomado como evidência. Nas pistas exploratórias o resultado foi misto: o 5×2 teve menor mediana na Wang 2.2 e o 5×3 na pista de treino. Os modelos 5×3 tiveram menor intervalo interquartil nas pistas de teste e menor erro no quartil de menor velocidade (mediana de 1,03° contra 1,37°); no quartil de maior velocidade, porém, o ângulo pós-calculado dos modelos 5×2 foi melhor (0,68° contra 0,87°). Esses são resultados descritivos, fora do plano de análise.

Quanto à segunda pergunta, a terceira saída tem custo: a estimativa de $v_x$ piorou de forma significativa, com efeito grande, em duas das três pistas de teste; a de $v_y$ não se alterou de forma significativa; o custo de cada treinamento aumentou cerca de 70\%; e 2 dos 30 treinamentos 5×3 foram reprovados, contra nenhum dos 5×2.

Quanto à terceira pergunta, a saída evoluída e o ângulo calculado a partir das velocidades do próprio modelo 5×3 são praticamente equivalentes. Como a diferença em $v_y$ não foi significativa, o mecanismo do ganho observado na Wang 3.2 não foi identificado.

As equações de $\beta$ recuperam parcialmente a estrutura cinemática $\dot\beta \approx a_y/v_x - \dot\psi$: um termo de persistência próximo de 1 e um termo linear em $a_y$ aparecem na quase totalidade dos modelos, mas o termo linear em $\dot\psi$ aparece em apenas 8 dos 28. O estimador cinemático simples derivou na pista de treino, a mais longa e com trechos de baixa velocidade. Um modelo linear com os mesmos regressores e sem realimentação teve RMSE de $\beta$ de 1,8° a 4,0° nas pistas de teste, maior do que o dos modelos MGGP, e o linear com realimentação divergiu; ainda assim, não se separou a contribuição dos termos não lineares do MGGP.

Com a evidência disponível, não é possível afirmar que a terceira saída melhora o ângulo de deriva de forma geral. Se o ângulo de deriva é a única grandeza de interesse e o custo adicional é aceitável, o modelo 5×3 só é uma opção se a referência original estiver correta, pois com a referência corrigida o 5×2 foi melhor; se as velocidades também são usadas, o modelo 5×2 com $\beta$ pós-calculado preserva a qualidade de $v_x$ com menor custo de treinamento.

As principais limitações, discutidas na \autoref{sec:limitacoes}, são a base de dados e a configuração únicas, o deslocamento da velocidade lateral de referência nas pistas de teste e o número reduzido de pistas. Como trabalhos futuros, sugere-se: (i) esclarecer a origem desse deslocamento, com a equipe que coletou os dados, para definir qual referência é a correta; (ii) incluir uma linha de base linear com realimentação estável, por exemplo ajustada por erro de simulação; (iii) avaliar a sensibilidade à unidade de $\beta$ e à normalização da aptidão; (iv) incluir o operador de divisão protegida, para que o MGGP represente diretamente $v_y/v_x$ e $a_y/v_x$; (v) repetir a comparação com a configuração T19, de menor custo, e com mais pistas e veículos; (vi) investigar a poda estrutural baseada no \textit{Forward-Regression Orthogonal Estimator} (FROE), que a IC não conseguiu usar com a arquitetura MIMO; e (vii) combinar o estimador cinemático com o sensor virtual, usando o MGGP apenas para corrigir a deriva da integração.
